% !PRINTED 184
\weberpage{184}
%<<CONT
\weq{1}{\begin{aligned}
\varphi(x) &= Y_1^2 + Y_2^2 + \cdots Y_\nu^2 - Y_1'^2 - Y_2'^2 - \cdots - Y_{\nu'}'^2\\
&= Z_1^2 + Z_2^2 + \cdots Z_\mu^2 - Z_1'^2 - Z_2'^2 - \cdots - Z_{\mu'}'^2
\end{aligned}}
be two decompositions of $\varphi(x)$ into squares. Thus let $Y_1, Y_2 \ldots Y_\nu, Y_1', Y_2' \ldots Y_{\nu'}'$ be homogeneous linear functions of $x$ between which no linear relation with constant coefficients holds, so that $\nu + \nu' \leqq m$; the same assumptions are made about the functions $Z_1, Z_2 \ldots Z_\mu, Z_1', Z_2' \ldots Z_{\mu'}'$, so that $\mu + \mu' \leqq m$ as well.

Suppose that
\[
\nu < \mu,
\]
then we can subject the variables $x$ to the linear equations
\weq{2}{\begin{aligned}
Y_1 &= 0, \quad Y_2 = 0 \ldots Y_\nu = 0\\
Z_1' &= 0, \quad Z_2' = 0 \ldots Z_{\mu'}' = 0
\end{aligned}}
whose number $\nu + \mu'$ is smaller than $m$.

Now if all the functions $Z_1, Z_2 \ldots Z_\mu$ were linearly dependent on the functions $Y_1, Y_2 \ldots Y_\nu, Z_1', Z_2' \ldots Z_{\mu'}'$, then from these $\mu$ equations one could, by eliminating the $\nu$ variables $Y_1, Y_2 \ldots Y_\nu$, derive a linear equation between the $Z_1, Z_2 \ldots Z_\mu, Z_1' \ldots Z_{\mu'}'$, which by hypothesis is \gesperrt{not} to hold. Hence the vanishing of all the $Z_1, Z_2 \ldots Z_\mu$ is not a necessary consequence of the equations (2), and we can add to the equations (2) one further equation that is not homogeneous, say
\[
Z_1 = 1,
\]
then for the $m$ unknowns $x$ we have a system of $m$ or fewer equations which are independent of one another and therefore do not contradict one another.

But then the second representation (1) gives a positive value for $\varphi(x)$, while the first gives a negative or vanishing value, and in this lies a contradiction. From this it follows that
\[
\nu \geqq \mu,
\]
and since one concludes in the same way
\[
\mu \geqq \nu,
\]
only $\nu = \mu$ remains. In the same way one can prove that $\nu' = \mu'$, whereby our theorem is completely proved.
